% 数学ノート集・共通テンプレート v2（積分論2026年度第1部の版面を参考）
% A5判。章内で定義・定理・例・注意等の番号を共通化。
% コンパイル： platex -interaction=nonstopmode math_note_template_v2.tex (2回)
%              dvipdfmx math_note_template_v2.dvi
\documentclass[a5paper,10pt,onecolumn]{jsarticle}
\usepackage[dvipdfmx,top=22mm,bottom=20mm,inner=19mm,outer=18mm,headheight=13pt,headsep=6mm]{geometry}
\usepackage{amsmath,amssymb,amsthm,mathtools}
\usepackage{enumitem,fancyhdr,titlesec}
\usepackage[dvipdfmx,bookmarks=true,colorlinks=false]{hyperref}
\usepackage{pxjahyper}

% ----- 差し替える文書情報 -----
\newcommand{\NoteTitle}{コンパクト性とTychonoffの定理}
\newcommand{\NoteSeries}{数学ノート集}

% ----- 文字・組版 -----
% 参考PDFの質素な明朝＋Computer Modernの印象を引き継ぎながら，
% A5判で窮屈にならないよう本文の行間を少し確保する．
\linespread{1.13}
\setlength{\parindent}{1zw}
\setlength{\parskip}{0pt}
\setlength{\abovedisplayskip}{6pt plus 2pt minus 2pt}
\setlength{\belowdisplayskip}{6pt plus 2pt minus 2pt}
\setlength{\abovedisplayshortskip}{3pt plus 2pt minus 1pt}
\setlength{\belowdisplayshortskip}{4pt plus 2pt minus 1pt}
\setlist[enumerate]{topsep=3pt,itemsep=1pt,parsep=0pt,leftmargin=2.4zw}
\titleformat{\section}{\normalfont\normalsize\bfseries}{\thesection}{0.8zw}{}
\titlespacing*{\section}{0pt}{2.0ex plus .4ex}{.9ex}
\titleformat{\subsection}{\normalfont\normalsize\bfseries}{\thesubsection}{0.7zw}{}
\titlespacing*{\subsection}{0pt}{1.5ex plus .3ex}{.6ex}
\raggedbottom

% ----- 柱とページ番号：ページ上部、罫線なし -----
\fancypagestyle{mathnote}{%
  \fancyhf{}%
  \fancyhead[L]{\footnotesize \NoteSeries\ （\NoteTitle）}%
  \fancyhead[R]{\footnotesize \thepage}%
  \renewcommand{\headrulewidth}{0pt}%
  \renewcommand{\footrulewidth}{0pt}%
}
\pagestyle{mathnote}

% ----- 節内通し番号の定義・定理群 -----
\newtheoremstyle{mathnotetheorem}
 {4.5pt}{4.5pt}{\normalfont}{}%
 {\bfseries}{}{0.65zw}%
 {\thmname{#1}\thmnumber{ #2}\thmnote{（#3）}}
\theoremstyle{mathnotetheorem}
\newtheorem{theorem}{定理}[section]
\newtheorem{definition}[theorem]{定義}
\newtheorem{lemma}[theorem]{補題}
\newtheorem{proposition}[theorem]{命題}
\newtheorem{corollary}[theorem]{系}
\newtheorem{example}[theorem]{例}
\newtheorem{remark}[theorem]{注意}
\renewcommand{\proofname}{証明}
\renewcommand{\qedsymbol}{$\square$}
\makeatletter
\renewenvironment{proof}[1][\proofname]{%
  \par\pushQED{\qed}\normalfont\topsep4\p@\@plus2\p@\relax
  \trivlist\item[\hskip\labelsep\bfseries #1]\ignorespaces
}{\popQED\endtrivlist\@endpefalse}
\makeatother

% ----- 使うノートで適宜拡張する記号 -----
\newcommand{\R}{\mathbb{R}}
\newcommand{\norm}[1]{\lVert #1\rVert}
\newcommand{\abs}[1]{\lvert #1\rvert}

\begin{document}
\section*{読了後にできること}
本稿を読むための前提は，開被覆，積位相，有限交叉性．
本稿の具体的な到達課題は，コンパクト性の証明法を比較し，積空間での有限交叉性を使う．


\section{コンパクト性と有限交叉性}
\begin{definition}
位相空間 $X$ がコンパクトであるとは，
$X$ の任意の開被覆が有限部分被覆を持つことをいう．
\end{definition}
\begin{proposition}[閉集合の有限交叉性]\label{prop:fip}
$X$ がコンパクトであることと，
閉集合族 $(F_i)$ の任意の有限部分族の共通部分が
非空なら全体の共通部分も非空となることは同値である．
\end{proposition}
\begin{proof}
$\bigcap_iF_i=\varnothing$ と
$\bigcup_i(X\setminus F_i)=X$ はDe Morganの法則で同値である．
後者に有限部分被覆があることは
前者が有限部分族について既に成立することを意味する．
従って対偶を取ればよい．
\end{proof}
\section{フィルターと超フィルター}
\begin{definition}
$\mathcal F\subset\mathcal P(X)$ がフィルターであるとは，
$\varnothing\notin\mathcal F$，
上位集合に関して閉じ，有限交叉に関して閉じることをいう．
極大フィルターを超フィルターという．
$\mathcal F$ が $x$ に収束するとは，
$x$ の近傍全てが $\mathcal F$ に含まれることをいう．
\end{definition}
\begin{lemma}[超フィルターへの拡張]\label{lem:ultra}
集合 $X$ 上の任意のフィルター $\mathcal F$ は
ある超フィルターに含まれる．
\end{lemma}
\begin{proof}
$\mathcal F$ を含むフィルター全体を包含順序で考える．
鎖の和集合は，任意の二つの元が鎖の一方のフィルターに
共に含まれるため，有限交叉について閉じる．
空集合を含まず上位集合についても閉じる．
従って鎖は上界を持ち，Zornの補題から極大元が存在する．
\end{proof}
\begin{lemma}[超フィルターの二値性]\label{lem:dichotomy}
超フィルター $\mathcal U$ と $A\subset X$ に対して
$A\in\mathcal U$ または $X\setminus A\in\mathcal U$ の
ちょうど一方が成立する．
\end{lemma}
\begin{proof}
両方含めば空集合を含んでしまう．
$A\notin\mathcal U$ なら，
$\mathcal U\cup\{A\}$ が生成する族は
超フィルターの極大性により真のフィルターでない．
従ってある $B\in\mathcal U$ に対して
$A\cap B=\varnothing$ となる．
これより $B\subset X\setminus A$ なので
上方閉性から $X\setminus A\in\mathcal U$ である．
\end{proof}
\begin{theorem}[コンパクト性の超フィルター判定]\label{thm:ultra}
$X$ がコンパクトであることと，
$X$ 上の全ての超フィルターが収束することは同値である．
\end{theorem}
\begin{proof}
$X$ をコンパクトとし，$\mathcal U$ を超フィルターとする．
$\{\,\overline A:A\in\mathcal U\,\}$ は有限交叉性を持つ：
有限個の $A_i\in\mathcal U$ の共通部分は空でないからである．
命題\ref{prop:fip}から
$x\in\bigcap_{A\in\mathcal U}\overline A$ を選ぶ．
もし $x$ の開近傍 $V$ が $\mathcal U$ に属さなければ，
補題\ref{lem:dichotomy}から $X\setminus V\in\mathcal U$．
だが $x\in\overline{X\setminus V}$ は
$V$ が開であることと矛盾する．
従って $\mathcal U\to x$ である．
逆に $X$ がコンパクトでなければ，
有限部分族が全体を覆わない開被覆 $(V_i)$ を取れる．
閉集合 $(X\setminus V_i)$ は有限交叉性を持つので
それらを含むフィルターを作り，
補題\ref{lem:ultra}で超フィルター $\mathcal U$ に延長する．
これが $x\in X$ に収束するとすれば
$x\in V_i$ となる $i$ を取り，
$V_i\in\mathcal U$ と $X\setminus V_i\in\mathcal U$
が同時に成立して矛盾する．
\end{proof}
\section{Tychonoffの定理}
\begin{theorem}[Tychonoff]\label{thm:tych}
コンパクト位相空間族 $(X_i)_{i\in I}$ の積
$\prod_{i\in I}X_i$ は積位相についてコンパクトである．
\end{theorem}
\begin{proof}
空の積は一点空間である．ある $X_i$ が空なら
積も空なのでコンパクトである．
以下全ての $X_i$ が非空の場合を考える．
積の上の超フィルター $\mathcal U$ を取る．
各射影 $\pi_i$ による像のフィルター
\[
 (\pi_i)_*\mathcal U=
 \{B\subset X_i:\pi_i^{-1}(B)\in\mathcal U\}
\]
は超フィルターである．
$X_i$ はコンパクトなので
定理\ref{thm:ultra}から
$(\pi_i)_*\mathcal U\to x_i$ となる
$x_i\in X_i$ が存在する．
選択公理により $x=(x_i)_{i\in I}$ を取る．
$x$ の積位相の基本近傍は
有限個の添字 $F$ だけで
$\pi_i(y)\in V_i$ を要求する形である．
各 $V_i$ が像フィルターに属すれば
$\pi_i^{-1}(V_i)\in\mathcal U$ であり，
有限交叉も $\mathcal U$ に属する．
従って $\mathcal U\to x$ である．
再び定理\ref{thm:ultra}より積空間はコンパクトである．
\end{proof}
\section{標準的な帰結}
\begin{proposition}[連続像・閉部分集合]
コンパクト空間の連続像はコンパクトである．
コンパクト空間の閉部分集合はコンパクトである．
\end{proposition}
\begin{proof}
連続写像 $f:X\to Y$ と $f(X)$ の開被覆を取る．
その逆像は $X$ の開被覆なので有限部分被覆があり，
その像は $f(X)$ の有限被覆を与える．
閉部分集合 $F\subset X$ の開被覆を取れば，
被覆に $X\setminus F$ を一つ加えて $X$ の開被覆になる．
有限部分被覆から $X\setminus F$ を除けば
$F$ の有限部分被覆を得る．
\end{proof}
\begin{proposition}[Hausdorff空間のコンパクト集合]
$Y$ がHausdorffなら任意のコンパクト部分集合 $K\subset Y$ は閉である．
\end{proposition}
\begin{proof}
$x\notin K$ に対し，
各 $y\in K$ に対して $x,y$ の互いに素な開近傍
$U_y,V_y$ を取る．
有限個の $V_{y_j}$ が $K$ を覆う．
$U=\bigcap_jU_{y_j}$ は $x$ の開近傍であり，
$K$ と交わらない．
従って $Y\setminus K$ は開である．
\end{proof}
\begin{corollary}[コンパクトからHausdorffへの連続全単射]
$X$ をコンパクト，$Y$ をHausdorffとする．
連続全単射 $f:X\to Y$ は同相写像である．
\end{corollary}
\begin{proof}
$X$ の閉集合 $F$ はコンパクトなので
$f(F)$ はコンパクト，従って $Y$ で閉である．
$f$ は閉写像であるから逆写像は連続となる．
\end{proof}

\section{到達課題の完全解答}
\noindent\textbf{積空間のコンパクト性を使った具体的な存在証明．}
\begin{proposition}[Cantorの積空間]
$X=\{0,1\}^{\mathbb N}$ に積位相を入れると
$X$ はコンパクトであり，各点を含む基本開集合は
有限個の座標の指定だけで与えられる．
\end{proposition}
\begin{proof}
各因子 $\{0,1\}$ は有限離散空間だからコンパクトである．
Tychonoffの定理を適用して積 $X$ はコンパクトとなる．
積位相の定義により基本開集合は，有限個の添字
$i_1,\ldots,i_r$ の座標を所定の値に固定する集合である．
\end{proof}
\begin{proposition}[有限条件の整合性から無限列を構成する]
$F_n\subset X$ を閉集合で，
$F_1\supset F_2\supset\cdots$ とする．
各 $F_n$ が非空なら $\bigcap_{n\ge1}F_n\ne\varnothing$．
\end{proposition}
\begin{proof}
任意の有限個の $F_n$ の交叉は，
添字が最大の集合そのものだから非空である．
コンパクト空間での閉集合の有限交叉性により
全体の交叉は非空である．
\end{proof}
\noindent\textbf{到達確認．}
$F_n=\{x\in X:x_1=\cdots=x_n=0\}$ の交叉を求めよ．
\par\noindent\textbf{解答．}
各座標が零である条件がすべて課されるので
$\bigcap_nF_n=\{(0,0,\ldots)\}$ である．
\end{document}
